%% Не менять - Do not modify
%%\input{my_folder/summary_settings} 
\chapter*[Count-me]{Реферат} % * - не нумеруем
\thispagestyle{empty}% удаляем параметры страницы
%\setcounter{sumPageFirst}{\value{page}}
%sumPageFirst \arabic{sumPageFirst}
%
%
%% Возможность проверить другие значения счетчиков - debugging
%\ref*{TotPages}~с.,
%\formbytotal{mytotalfigures}{рисун}{ок}{ка}{ков},
%\formbytotal{mytotaltables}{таблиц}{у}{ы}{},
%There are \TotalValue{mytotalfigures} figures in this document
%There are \TotalValue{mytotalfiguresInApp} figuresINAPP in this document
%There are \TotalValue{mytotaltables} tables in this document
%There are \TotalValue{mytotaltablesInApp} figuresINAPP in this document
%There are \TotalValue{myappendices} appendix chapters in this document
%\total{citenum}~библ. наименований.



%% Для того, чтобы значения счетчиков корректно отобразились, необходимо скомпилировать файл 2-3 раза
%На \total{mypages}~c.,  
%\formbytotal{myfigures}{рисун}{ок}{ка}{ков},
%\formbytotal{mytables}{таблиц}{у}{ы}{},
%\formbytotal{myappendices}{приложен}{ие}{ия}{ий}%.  

%% 2025!!!
%% Для того, чтобы значения счетчиков корректно отобразились, необходимо скомпилировать файл 2-3 раза
На \total{mypages}~c.,  
\total{myfigures}~{рис.},
\total{mytables}~{табл.},
\total{myappendices}~{прил.}%.



%\noindent
{Ключевые слова: \keywordsRu} % Ключевые слова из renames.tex

Тема выпускной квалификационной работы: <<\thesisTitle>>. % изменения 2025!


%\abstractRu % Аннотация из renames.tex
В работе разработана архитектура модели глубокого обучения для прогнозирования 
фенотипических признаков сельскохозяйственных культур на основе SNP-данных. 
Исследование проведено на трёх культурах: нут, рис и рапс. Рассматривались 
масса тысячи семян, высота растения, время цветения, содержание масла и 
глюкозинолатов.

Предложенная модель включает четыре блока: генетического кодирования 
с механизмом внимания на уровне SNP, трансформерный для моделирования 
межгенных взаимодействий, прогнозирования фенотипов и реконструкционный 
для регуляризации скрытых представлений. Для обучения применена многозадачная 
функция потерь с адаптивным взвешиванием признаков.

Проведена серия вычислительных экспериментов: сравнительный анализ методов 
аугментации данных, исследование стратегий отбора генов и оценка влияния 
архитектурных модификаций на точность прогнозирования. Для интерпретации 
модели использованы значения Шепли, позволившие выявить значимые генные 
взаимодействия.

Наибольшая точность прогнозирования достигнута для времени цветения рапса 
(коэффициент корреляции 0.809) и массы тысячи семян нута (0.881). 
Эффективность подхода подтверждена на данных различных культур, что 
демонстрирует его обобщающую способность. Разработанная модель может 
быть применена в селекционных программах для идентификации ключевых 
генетических маркеров, связанных с продуктивностью растений.
%\abstractRu % УДАЛИТЬ. Повтор иллюстрации переноса текста на вторую страницу



\printTheAbstract % не удалять


On~\total{mypages}~pages, 
\total{myfigures}~figures, 
\total{mytables}~tables,
\total{myappendices}~appendices%.

%\noindent
{Keywords: \keywordsEn} % Ключевые слова из renames.tex 
	
The subject of the graduate qualification work is <<\thesisTitleEn>>.
	
	
A deep learning architecture for predicting phenotypic traits of agricultural crops from SNP data is developed. The study covers three crops: chickpea, rice and rapeseed. Target traits include thousand seed weight, plant height, flowering time, oil content and glucosinolate content.

The model comprises four blocks: genetic encoding with SNP-level attention, 
transformer for intergenic interactions, phenotype prediction and reconstruction 
for regularisation. A multi-task loss with adaptive weighting is employed.

Computational experiments include augmentation methods comparison, gene 
selection strategies and architectural modifications. Shapley values reveal 
significant gene interactions.

Top accuracy: rapeseed flowering time (0.809) and chickpea thousand seed 
weight (0.881). The approach generalises across crops and can be applied 
in breeding programmes.
%\abstractEn % Аннотация из renames.tex

%\abstractEn % УДАЛИТЬ. Повтор для иллюстрации переноса текста на вторую страницу
	


%% Не менять - Do not modify
\thispagestyle{empty}
%\setcounter{sumPageLast}{\value{page}} %сохранили номер последней страницы Задания
%\setcounter{sumPages}{\value{sumPageLast}-\value{sumPageFirst}}
%sumPageLast \arabic{sumPageLast}
%
%sumPages \arabic{sumPages}
%\restoregeometry % восстанавливаем настройки страницы
%\input{my_folder/summary_settings_restore}	% настройки - конец